\begin{tabularx}{\textwidth}{L{32mm}X}
\toprule
\textbf{Tulajdonság} & \textbf{Jelentés} \\
\midrule
Indexbejegyzés & Kulcsértékből és rekordra vagy adatblokkra mutató pointerből áll. \\
Rendezettség & A kulcsok rendezettek, ezért a keresés nem teljes táblabejárással történik. \\
B-fa csomópont & Több kulcsot és több gyermekpointert tartalmaz; egy csomópont gyakran lemezblokknak felel meg. \\
Kiegyensúlyozottság & A levelek azonos mélységben vannak, ezért a keresés útja rövid és kiszámítható. \\
Beszúrás & A megfelelő levélbe történik; telített csomópont esetén hasítás, középső kulcs felvitele és pointerigazítás szükséges. \\
Törlés & Hiányos csomópont esetén kölcsönzés vagy összevonás történhet, hogy a fa kiegyensúlyozott maradjon. \\
Költség & Tipikusan \(O(\log n)\) blokkelérés; a gyakorlatban a fa nagy elágazási foka miatt kevés szint elegendő. \\
Ár & A lekérdezések gyorsulhatnak, de az \code{INSERT}, \code{UPDATE}, \code{DELETE} műveleteknél az indexeket is frissíteni kell. \\
\bottomrule
\end{tabularx}
